\section{显式 3D 表示}

3D 物体的几何表示可以分为两大类：\textbf{显式表示}（Explicit）和\textbf{隐式表示}（Implicit）。显式表示直接给出物体表面或内部的具体坐标信息。

\begin{figure}[H]
\centering
\includegraphics[width=0.95\textwidth]{slides-images/slide-003.jpg}
\caption{3D 表示的两大分类：显式（点云、网格、参数化曲面）与隐式（水平集、代数曲面、距离函数）\protect\footnotemark}
\end{figure}
\footnotetext{Slides 第3页。}

\subsection{点云（Point Cloud）}

点云是最简单的 3D 表示，由一组 3D 点的坐标构成，存储为 $3 \times N$ 矩阵（$xyz$ 坐标 $\times$ 点数）。

\begin{figure}[H]
\centering
\includegraphics[width=0.95\textwidth]{slides-images/slide-004.jpg}
\caption{点云表示：每个点有 $xyz$ 坐标，可选附加法向量信息（Surfels）\protect\footnotemark}
\end{figure}
\footnotetext{Slides 第4页。}

\textbf{点云的优势}：
\begin{itemize}
    \item 是 3D 传感器（深度相机、3D 扫描仪、iPhone AR Kit 等）的\textbf{原始输出格式}
    \item 非常灵活，不受拓扑约束，可以表示任意形状
    \item 适合处理大规模多样化数据集
\end{itemize}

\textbf{点云的局限}：
\begin{itemize}
    \item 采样可能不均匀（兔子头部点多，尾部点少）
    \item 没有连接信息，无法直接进行细分或简化操作
    \item \textbf{缺乏拓扑信息}——仅凭点无法区分圆环（torus）和球体
    \item 不支持平滑渲染
\end{itemize}

\begin{figure}[H]
\centering
\includegraphics[width=0.95\textwidth]{slides-images/slide-005.jpg}
\caption{点云的来源与应用：3D 传感器扫描得到点云，需后处理融合为完整模型\protect\footnotemark}
\end{figure}
\footnotetext{Slides 第5页。}

\begin{warningbox}{点云缺乏拓扑信息}
给定一组点，无法判断物体是实心的还是有孔的。例如，圆环和球体的点云在某些采样下可能看起来相似。这是点云表示的根本局限。
\end{warningbox}

\subsection{多边形网格（Polygon Mesh）}

为了补充点云缺乏的连接信息，自然引入了\textbf{多边形网格}——不仅存储顶点，还存储面（面通常是三角形或四边形）。

\begin{figure}[H]
\centering
\includegraphics[width=0.95\textwidth]{slides-images/slide-006.jpg}
\caption{多边形网格：顶点 + 面（连接关系），是图形引擎和游戏中最广泛使用的表示\protect\footnotemark}
\end{figure}
\footnotetext{Slides 第6页。}

网格是图形引擎和游戏中\textbf{最广泛使用}的 3D 表示。复杂的网格可以达到数千万三角形（如 5600 万三角形的 David 雕塑），甚至 Google Earth 使用\textbf{万亿级}三角形来表示地球表面。

\begin{figure}[H]
\centering
\includegraphics[width=0.95\textwidth]{slides-images/slide-007.jpg}
\caption{高精度网格示例：David 雕塑（5600万三角形）和大规模城市模型\protect\footnotemark}
\end{figure}
\footnotetext{Slides 第7页。}

网格支持丰富的操作：
\begin{itemize}
    \item \textbf{细分（Subdivision）}：增加三角形数量以捕捉更多细节
    \item \textbf{简化（Simplification）}：减少三角形数量以提高处理速度
    \item \textbf{正则化（Regularization）}：使三角形大小均匀、形状规则
\end{itemize}

\begin{warningbox}{网格与深度学习的兼容性挑战}
网格的不规则性（每个面可能有不同数量的顶点，分辨率可变）使其难以直接输入标准 CNN。这是 3D 深度学习起步较晚的重要原因之一——人们需要找到将不规则 3D 数据适配到神经网络的方法。
\end{warningbox}

\subsection{参数化表示（Parametric Representations）}

对于具有规则结构的物体（如圆、球、工业零件），可以用\textbf{参数函数}来表示其形状。

\begin{figure}[H]
\centering
\includegraphics[width=0.95\textwidth]{slides-images/slide-008.jpg}
\caption{参数化表示：用函数将低维参数空间映射到 3D 空间。圆可用 $(\cos t, \sin t)$ 表示，球可用两个参数 $(u, v)$ 映射\protect\footnotemark}
\end{figure}
\footnotetext{Slides 第8页。}

$$\text{圆：} \quad \mathbf{p}(t) = (\cos t, \sin t), \quad t \in [0, 2\pi]$$
$$\text{球：} \quad \mathbf{p}(u, v) = (\sin u \cos v, \sin u \sin v, \cos u)$$

更复杂的参数化表示包括 \textbf{B\'{e}zier 曲线/曲面}和 \textbf{B-Spline}，通过少量控制点就能表示光滑的自由曲面。

\begin{importantbox}{显式表示的核心特点}
显式表示\textbf{易于采样}——给定参数值，直接计算得到物体表面上的点。但\textbf{难以判断}某个任意点是在物体内部还是外部。
\end{importantbox}

\subsection{本章小结}

\begin{itemize}
    \item 点云：最简单，是传感器的原始格式，但缺乏拓扑信息
    \item 网格：最广泛使用，支持丰富操作，但不规则性使其难以直接与深度学习集成
    \item 参数化曲面：用函数紧凑地表示规则形状，支持细分操作
    \item 显式表示共同的优势是\textbf{采样容易}，共同的劣势是\textbf{内外判断困难}
\end{itemize}

%% ========== 隐式表示 ==========

